\documentclass{leveldoc}
\DocUses{foundations,deflection_plates}
\DocFingerprint{bf08ca5fa8915aaf}

\begin{document}

\section{What the level asks}\label{sec:asks}

\leveldescription

One positive ion is launched along $+x$ from the left edge. Above its line a
strong positive charge is fixed, and between the two lies a bar of glass, an
insulator of relative permittivity $\varepsilon$. The glass carries no free
charge, and none can flow onto it, but in the charge's field its molecules
polarize, and a bound charge appears on its surfaces. The player may add up to
three charges, of either sign, from the listed magnitudes, inside the marked
region below the launch line. There are no magnets and no electrodes.

The ion must enter the detector. The shot has no acceptance window, no kinetic
window and no gate. A flight ends when the ion touches a charge or the glass,
leaves the world, or enters the detector (\xsecref{foundations}{sec:board}).
The glass's obstacle is its box plus a contact shell of thickness
$\val{contact}$, as for an electrode. A fixed charge has radius $\val{Rcharge}$,
and the ion's radius is $0$.

With nothing placed, the fixed charge throws the ion down and out through the
bottom of the world. The reference adds one negative charge low on the board,
ahead of the fixed charge's line: the ion swings round it and comes up into the
detector's left face. Removed from the board, the glass changes the ending: the
same flight then strikes the placed charge. The glass does not screen the
ion's line from the fixed charge. The bound charge on its far face, nearer the
ion, pushes along with the fixed charge, and the push on the launch line is
stronger than without the glass; only inside the glass is the field weaker.
This document states the bound-charge problem, gives the infinite slab as its
analytic limit, solves it on rectangles, and integrates the flights by the
classical Runge--Kutta method.

\section{Setup}\label{sec:setup}

\setuptable

The speed of light is $c = \val{cLight}$. The ion has charge $q = \val{qPart}$,
mass $m = \val{mPart}$ and radius $0$, and it is launched with kinetic energy
$T_0 = \val{T0}$ from $(\val{xLaunch}, \val{yLaunch})$, straight along $+x$.
Then $\gamma_0 = \val{gamma0}$, $\abs{\vect{p}_0} = \val{p0}$ and
$v_0 = \val{v0} = \val{vOverC}\,c$ (\xeqref{foundations}{eq:launch}).

The fixed charge is $Q = \val{qFixed}$ at $(\val{xCharge}, \val{yCharge})$.
The glass occupies $\val{xGlass0} \le x \le \val{xGlass1}$ and
$\val{yGlass0} \le y \le \val{yGlass1}$: length $\val{glassL}$ along $x$,
thickness $\val{glassT}$ along $y$, and height $\val{glassH}$ along $z$,
symmetric about the plane $z = 0$ of the flight. Its permittivity is
$\varepsilon = \val{eps}$, and we write
\begin{equation}\label{eq:lambda}
  \lambda = \frac{\varepsilon - 1}{\varepsilon + 1} = \val{lam}.
\end{equation}
The reference charge is $Q_\mathrm{r} = \val{qPlay}$ at $(\val{xPlay}, \val{yPlay})$.
The detector is $\val{xDetLo} \le x \le \val{xDetHi}$,
$\val{yDetLo} \le y \le \val{yDetHi}$. The world is $-1 \le x \le \val{xWorld}$,
$-1 \le y \le \val{yWorld}$. The reported time step is $h = \val{H}$; the
reported mesh uses rectangles of edge $\val{panelFine}$, $\val{nPanels}$ of
them on the upper half of the glass.

\begin{figure}[ht]
\centering
\begin{tikzpicture}
\begin{axis}[axis equal image, width=0.92\linewidth,
  xmin=-1.5, xmax=32.5, ymin=-2, ymax=22.5,
  xlabel={$x$}, ylabel={$y$}, tick label style={font=\small}]
  \draw[black!40, dashed] (-1,-1) rectangle (\rawval{xWorld},\rawval{yWorld});
  \draw[black!25] (0,0) rectangle (30,20);
  \fill[green!50!black, fill opacity=0.18]
    (\rawval{xDetLo},\rawval{yDetLo}) rectangle (\rawval{xDetHi},\rawval{yDetHi});
  \draw[green!40!black]
    (\rawval{xDetLo},\rawval{yDetLo}) rectangle (\rawval{xDetHi},\rawval{yDetHi});
  \fill[cyan!40!white]
    (\rawval{xGlass0},\rawval{yGlass0}) rectangle (\rawval{xGlass1},\rawval{yGlass1});
  \draw[cyan!50!black]
    (\rawval{xGlass0},\rawval{yGlass0}) rectangle (\rawval{xGlass1},\rawval{yGlass1});
  \fill[red!70!black] (\rawval{xCharge},\rawval{yCharge}) circle[radius=\rawval{Rcharge}];
  \fill[blue!70!black] (\rawval{xPlay},\rawval{yPlay}) circle[radius=\rawval{Rcharge}];
  \addplot[orange!80!black] table {pathbare.dat};
  \addplot[very thick, blue!60!black] table {pathref.dat};
  \draw[-{Stealth[length=5pt]}, thick] (\rawval{xLaunch},\rawval{yLaunch}) -- (\rawval{arrowX},\rawval{arrowY});
  \fill (\rawval{xLaunch},\rawval{yLaunch}) circle[radius=0.08];
\end{axis}
\end{tikzpicture}
\caption{The glass (cyan), the fixed charge (red, above it) and the reference
charge (blue, low on the right); the green box is the detector. Blue is the
reference flight: it dips towards the placed charge, swings round it and comes
up into the detector's left face. Orange is the same launch with nothing
placed: the fixed charge throws it out through the bottom of the world.}
\label{fig:board}
\end{figure}

\section{Principles and model}\label{sec:model}

\begin{principles}
  \item The field of a point charge \xeqref{foundations}{eq:coulomb}, and the
    force on a non-polarizable particle \xeqref{foundations}{eq:sphereforce}.
    There is no magnetic field.
  \item The relativistic equation of motion \xeqref{foundations}{eq:eom}, the
    launch \xeqref{foundations}{eq:launch}, and energy in a static field
    \xeqref{foundations}{eq:energy}. The motion stays in the plane
    (\xsecref{foundations}{sec:plane}).
  \item A linear dielectric: its polarization is proportional to the field in
    it, so $\vect{D} = \varepsilon\vect{E}$ there, and with no free charge on
    its surface the normal component of $\vect{D}$ is continuous across it
    (\Jackson{§4.3, eq.~(4.40)}, in the Gaussian units of the Foundations). Its response is
    instantaneous on the flight's time scale (\physref{2.7}).
  \item The potential and field of a uniformly charged rectangle
    \xeqref{deflection_plates}{eq:rectangle}, with its primitive
    \xeqref{deflection_plates}{eq:primitive}.
  \item The board's rules (\xsecref{foundations}{sec:board}), the classical
    Runge--Kutta method and the cubic Hermite location of an event
    (\xsecref{foundations}{sec:numerics}), as in the deflection plates' flight
    (\xsecref{deflection_plates}{sec:flight}).
  \item The radiated power \xeqref{foundations}{eq:larmor}, summed by Simpson's
    rule \xeqref{foundations}{eq:simpson}. It is not put back into the force.
\end{principles}

The polarized glass acts on the ion through its bound charge only: a surface
density $\sigma$ on its faces (the volume density
$-\nabla\cdot\vect{P}$ vanishes in a uniform linear medium without free charge).
Off the glass the field is the Coulomb field of the fixed and placed charges
plus the field of $\sigma$ and of its mirror in the plane $z = 0$. The ion is
not among the charges that polarize the glass. The state $(\vect{x}, \vect{p})$
obeys \xeqref{foundations}{eq:eom} with $\vect{B} = \vect{0}$ and with that
$\vect{E}$, from the launch until the first ending event.

\begin{assumption}{A static, linear glass}\label{as:static}
The glass answers at once, with the same $\varepsilon$ at every field: no
dispersion, no losses, no breakdown, no charge collected from the beam. The
launch speed is $\val{vOverC}\,c$. The game makes the same assumptions
(\physref{2.7}), so this is the model being compared.
\end{assumption}

\begin{assumption}{No radiation reaction}\label{as:radiation}
The ion radiates \xeqref{foundations}{eq:larmor}, and that energy is not taken
from the flight. On the reference flight the radiated energy is a fraction
$\val{RefRadFrac}$ of $T_0$.
\end{assumption}

\begin{assumption}{No image of the ion}\label{as:image}
The bound charge the ion itself induces on the glass is left out: it would need
a new solve at every force evaluation. The game leaves it out for the same
reason (\physref{2.7}). A conductor's image force bounds it, $q^2/d^2$ with the
distance $d$ to the glass floored at $10^{-3}$ (a dielectric's image is the
conductor's times $\lambda < 1$), taken relative to $\max(\abs{q\vect{E}}, T_0)$.
Along the reference samples the ratio reaches $\val{RefImage}$; along the
empty board's samples it reaches $\val{BareImage}$. The engine asks for it to
stay below $10^{-10}$.
\end{assumption}

\begin{assumption}{Equal rectangles}\label{as:panels}
The bound charge is constant on equal rectangles of edge $\val{panelFine}$,
and the surface condition is held on average over each, at four Gauss points.
The engine holds the same condition on graded triangles (\physref{2.7}).
Replacing the edge by $\val{panelCoarse}$ moves an end state by
$\val{errMesh}$; that change is the discretization estimate of this mesh.
\end{assumption}

\begin{assumption}{The fixed step}\label{as:step}
The reported step is $h = \val{H}$. Halving it estimates an end state's error
by $\val{errStep}$ (\xsecref{foundations}{sec:numerics}). Ten times the panel
estimate $\val{errMesh}$ is above $10^{-9}$ and above ten times the step
estimate, so the tolerance on a state is $\val{tolFlight}$. Between samples
the ion moves at most $\val{stepReach}$, while its smallest sampled clearance
of the glass is $\val{RefDmin}$ on the reference flight and $\val{BareDmin}$
on the empty board.
\end{assumption}

\section{Derivation}\label{sec:derivation}

\subsection{The surface condition}\label{sec:condition}

Take a point of the glass's surface with outward normal $\uvec{n}$, and let
$\bar E_n$ be the normal field there of every charge except the surface
element itself: the fixed and placed charges, the rest of the bound charge,
and the mirror. A small patch of density $\sigma$ adds $+2\pi\sigma$ to the
normal field just outside and $-2\pi\sigma$ just inside, by Gauss's law on a
thin box round the patch with the field \xeqref{foundations}{eq:coulomb}, so
\[
  E_n^\mathrm{out} = \bar E_n + 2\pi\sigma, \qquad
  E_n^\mathrm{in} = \bar E_n - 2\pi\sigma .
\]
Outside is vacuum, inside $\vect{D} = \varepsilon\vect{E}$, and with no free
charge the normal $\vect{D}$ is continuous:
$\varepsilon(\bar E_n - 2\pi\sigma) = \bar E_n + 2\pi\sigma$. Solving for
$\sigma$,
\begin{equation}\label{eq:condition}
  \sigma = \frac{\lambda}{2\pi}\,\bar E_n ,
\end{equation}
with $\lambda$ of \eqref{eq:lambda}. Where the field leaves the glass
($\bar E_n > 0$) the bound charge is positive. The fixed charge is positive and
above the glass, so its field enters the near face (negative bound charge) and
leaves through the far face (positive bound charge), the face nearer the ion.

The glass is neutral, so its bound charge sums to zero:
\begin{equation}\label{eq:neutral}
  \oint \sigma \,\dd S = 0 .
\end{equation}
Equation \eqref{eq:condition} holds this only approximately once it is
discretized: on a closed body the uniform normal field is nearly a solution of
the homogeneous problem (the field of a net charge, which it would let in).
The discrete system therefore carries \eqref{eq:neutral} as a constraint with
one more unknown $\gamma$, a uniform normal field added on the whole surface,
\begin{equation}\label{eq:bordered}
  \sigma_i = \frac{\lambda}{2\pi}\Big(\bar E_{n,i} + \gamma\Big), \qquad
  \sum_i \sigma_i A_i = 0,
\end{equation}
where $i$ runs over the rectangles of \S\ref{sec:panels}, $A_i$ is a
rectangle's area, and the sum counts the upper half (the mirror doubles it).

\subsection{The infinite slab}\label{sec:slab}

An infinite slab of the glass's thickness $d = \val{slabD}$, with the charge
$Q$ a height $a = \val{slabA}$ above its near face, has a closed solution, the
limit of a bar that is long and tall compared with these distances. Put the
near face at $z' = 0$, the far face at $z' = -d$ and the charge at $z' = a$, on
the axis $\rho = 0$. In each region the potential is a superposition of the
solutions $J_0(k\rho)\,e^{\pm k z'}$ of Laplace's equation, and the charge's
own potential is $Q/\abs{\vect{r} - a\uvec{z}} = Q\int_0^\infty J_0(k\rho)\,
e^{-k\abs{z' - a}}\,\dd k$ (\Jackson{Ch.~3, cylindrical coordinates}). Write
\begin{align*}
  \phi_1 &= Q\!\int_0^\infty\! J_0(k\rho)\big[e^{-k\abs{z' - a}} + R\,e^{-k(z' + a)}\big]\dd k
    && (z' > 0),\\
  \phi_2 &= Q\!\int_0^\infty\! J_0(k\rho)\big[A\,e^{kz'} + B\,e^{-kz'}\big]\dd k
    && (-d < z' < 0),\\
  \phi_3 &= Q\!\int_0^\infty\! J_0(k\rho)\,C\,e^{kz'}\,\dd k
    && (z' < -d).
\end{align*}
At each face $\phi$ and $\varepsilon\,\partial\phi/\partial z'$ are continuous.
At $z' = -d$: $A e^{-kd} + B e^{kd} = C e^{-kd}$ and
$\varepsilon(A e^{-kd} - B e^{kd}) = C e^{-kd}$, so
$B = \lambda A\,e^{-2kd}$ and $C = (1 + \lambda)A$. At $z' = 0$:
$(1 + R)e^{-ka} = A + B$ and $(1 - R)e^{-ka} = \varepsilon(A - B)$; adding
them, $2e^{-ka} = A\,(1 + \varepsilon)\,(1 - \lambda^2 e^{-2kd})$. Hence
\[
  A = \frac{2\,e^{-ka}}{(1 + \varepsilon)(1 - \lambda^2 e^{-2kd})}, \qquad
  C = \frac{(1 - \lambda^2)\,e^{-ka}}{1 - \lambda^2 e^{-2kd}},
\]
using $2(1 + \lambda)/(1 + \varepsilon) = 4\varepsilon/(1 + \varepsilon)^2 =
1 - \lambda^2$. Expanding $1/(1 - \lambda^2 e^{-2kd}) = \sum_n \lambda^{2n}
e^{-2nkd}$ and integrating term by term, each term is a point charge's
potential again. Beyond the slab, a distance $h = \val{slabH}$ below the far
face, the field on the axis relative to the bare charge's, whose distance is
$r_0 = a + d + h = \val{rSample}$, is
\begin{equation}\label{eq:slab-out}
  \frac{E}{E^\mathrm{bare}} = (1 - \lambda^2) \sum_{n=0}^\infty
    \lambda^{2n}\Big(\frac{r_0}{r_0 + 2nd}\Big)^2 = \val{slabOut},
\end{equation}
a series of images at heights $a + 2nd$, the first carrying
$1 - \lambda^2 = \val{slabFirst}$ of $Q$. Inside, at the slab's middle plane,
a distance $z_\mathrm{c} = a + d/2 = \val{slabZc}$ from the charge, the
$A$ terms are charges at $a + 2nd$ and the $B$ terms charges below, at
$-(a + 2d + 2nd)$, and
\begin{equation}\label{eq:slab-in}
  \frac{E}{E^\mathrm{bare}} = \frac{2}{\varepsilon + 1} \sum_{n=0}^\infty
    \lambda^{2n}\bigg[\Big(\frac{z_\mathrm{c}}{z_\mathrm{c} + 2nd}\Big)^2
    - \lambda\Big(\frac{z_\mathrm{c}}{z_\mathrm{c} + d + 2nd}\Big)^2\bigg]
    = \val{slabIn}.
\end{equation}
The terms fall as $\lambda^{2n}$; sixty of them leave less than $10^{-26}$. So
even an endless slab leaves $\val{slabOut}$ of the field on the launch line,
not $1/\varepsilon = \val{invEps}$; inside, it keeps $\val{slabIn}$. A
dielectric is no screen: a conductor's surface charge cancels the field
beyond it, a dielectric's cannot, since by \eqref{eq:condition} it only follows
the field that crosses it.

The glass of this level is not a slab: it is $\val{glassH}$ high in $z$, about
its thickness, and its bound charge sits on two narrow strips facing the
charge and the ion. The field of such a strip pair beyond its far face is the
field of a dipole layer of finite width, which points along the dipoles below
it: the far face's positive charge, nearer the ion, outweighs the near face's
negative charge. The panels of \S\ref{sec:panels} give the numbers.

\subsection{The panels}\label{sec:panels}

The glass is a box. Only the upper half is meshed: the four side faces from
$z = 0$ up to half the height, and the top face. Each edge from $a$ to $b$ is
split into $n = \max(1, \operatorname{round}((b - a)/\val{panelFine}))$ equal
parts. The mirror of a rectangle through $z = 0$ carries the same density.
Rectangle $j$ at unit density has, at a point, the field
$-\nabla$\xeqref{deflection_plates}{eq:rectangle}, written $\vect{e}_j$; the
field of its mirror is $\vect{e}_j^\mathrm{m}$. With $\uvec{n}_i$ the outward
normal of rectangle $i$, the average of \eqref{eq:bordered} over rectangle $i$
at its four Gauss points $\vect{g}_{i,k}$ (the product rule of order three,
$\pm 1/\sqrt{3}$ of the half-edges) reads
\begin{equation}\label{eq:system}
  \sigma_i - \frac{\lambda}{2\pi}\sum_j \sigma_j K_{ij} - \frac{\lambda}{2\pi}\gamma
    = \frac{\lambda}{2\pi}\,\frac14\sum_k \vect{E}_\mathrm{ext}(\vect{g}_{i,k})\cdot\uvec{n}_i,
  \qquad
  K_{ij} = \frac14\sum_k \big(\vect{e}_j + \vect{e}_j^\mathrm{m}\big)(\vect{g}_{i,k})\cdot\uvec{n}_i ,
\end{equation}
where $\vect{E}_\mathrm{ext}$ is the Coulomb field of the fixed and placed
charges. In $K_{ii}$ the rectangle's own field is left out: its normal part at
its own points is $\pm 2\pi$, the jump already taken apart in
\S\ref{sec:condition}, and its average of the two sides is zero. Every
$K_{ij}$ projects on the receiving rectangle's normal $\uvec{n}_i$. With the
constraint \eqref{eq:neutral}, $\sum_i \sigma_i A_i = 0$, this is a dense
linear system for the $\val{nPanels}$ densities and $\gamma$, solved by LU
factorization. Its residual on the reported mesh is $\val{solveRes}$, and the
bound charge it leaves, both halves, is $\val{qNet}$ against
$Q = \val{qFixed}$.

The field at a point of the plane $z = 0$ is the Coulomb field of the charges
plus $\sum_j \sigma_j(\vect{e}_j + \vect{e}_j^\mathrm{m})$, whose normal
components in $z$ cancel there, so the ion stays in the plane. On the reported
mesh the near face's mean density is $\val{sigNear}$: negative, as
\eqref{eq:condition} says. The field against the bare charge's at a few
points, with nothing placed:
\begin{center}\small\val{fieldtab}\end{center}
Inside the glass the field falls to $\val{insideRel}$ of the bare charge's
(the slab's \eqref{eq:slab-in}: $\val{slabIn}$). On the launch line under the
charge it is $\val{screenRel}$ of it, larger, against the slab's
$\val{slabOut}$; under the glass's left end, $\val{endRel}$.

\subsection{The flight}\label{sec:flight}

There is no closed orbit. The state is advanced by \xeqref{foundations}{eq:rk4}
with the fixed step $h = \val{H}$, and the ending is located on the cubic
Hermite interpolant of each step, sampled at thirty-two fractions and refined
by sixty halvings, obstacles before the world and the world before the
detector, as in \xsecref{deflection_plates}{sec:flight}. The conserved energy
is
\begin{equation}\label{eq:flight-energy}
  W = (\gamma - 1)mc^2 + q\,\phi(\vect{x}),
\end{equation}
with $\phi$ the potential of the charges, the rectangles and their mirrors
(the ion is not in it). The step's error is $\abs{y_h - y_{2h}}/15$; the
panel estimate is the change in an end state when the edge doubles, not
divided. Both flights give the same outcome at $h$ and $2h$ and on both
meshes.

\section{The reference solution by hand}\label{sec:hand}

\begin{procedure}
  \item\label{st:launch} \emph{The launch} \xeqref{foundations}{eq:launch}.
    $\gamma_0 = \val{gamma0}$, $\abs{\vect{p}_0} = \val{p0}$, so
    $\vect{p}_0 = (\val{p0}, 0)$ and $v_0 = \val{v0}$.
  \item\label{st:slab} \emph{The limit.} With $\lambda = \val{lam}$,
    $a = \val{slabA}$, $d = \val{slabD}$, $h = \val{slabH}$, the series
    \eqref{eq:slab-out} gives $\val{slabOut}$ of the bare field on the launch
    line and \eqref{eq:slab-in} $\val{slabIn}$ inside. The march uses the
    panels, not the slab.
  \item\label{st:panels} \emph{The panels.} Build the upper-half rectangles of
    edge $\val{panelFine}$ as in \S\ref{sec:panels}: $\val{nPanels}$ of them,
    plus the mirrors. Solve \eqref{eq:system} with the fixed and the reference
    charge on its right side, and \eqref{eq:neutral}. The residual is
    $\val{solveRes}$.
  \item\label{st:march} \emph{The march.} Advance
    \xeqref{foundations}{eq:eom} by \xeqref{foundations}{eq:rk4} with step
    $\val{H}$ in that field, and locate the ending as in \S\ref{sec:flight}.
    The state every five time units:
    \begin{center}\small\val{states}\end{center}
  \item\label{st:end} \emph{The arrival.} The ion meets the \val{RefFace}
    face of the detector at $(\val{RefX}, \val{RefY})$, at $t = \val{RefT}$,
    with momentum $(\val{RefPx}, \val{RefPy})$, at $\val{RefAng}$ degrees above
    $+x$, after $\val{RefSteps}$ steps. Its lowest point is $y = \val{RefYmin}$,
    on its swing round the placed charge.
  \item\label{st:check} \emph{The checks.} Along the flight $W$ drifts by at
    most $\val{RefWdev}$. The step estimate of the end state is $\val{RefErr}$,
    the larger of the two flights' $\val{errStep}$; the panel estimate is
    $\val{errMesh}$. The sampled clearance of the glass is $\val{RefDmin}$,
    outside the shell $\val{contact}$, and the ion moves at most
    $\val{stepReach}$ in a step. The image-force ratio reaches $\val{RefImage}$.
  \item\label{st:rad} \emph{The radiation.} Simpson's rule on steps of $h/2$
    gives $\val{RefRad}$, a fraction $\val{RefRadFrac}$ of $T_0$; halving that
    step again changes it by the relative amount $\val{RefRadErr}$. It is our
    field's integral, not a row of the comparison.
\end{procedure}

\begin{result}[The reference flight]
The ion enters the detector through its \val{RefFace} face at
$(\val{RefX}, \val{RefY})$ at $t = \val{RefT}$, with momentum
$(\val{RefPx}, \val{RefPy})$. The level is solved.
\end{result}

\section{Comparison with the game}\label{sec:comparison}

\comparisontable

The table compares the reference's outcome, time and end point, and the same
for the empty board. The tolerance on a state, $\val{tolFlight}$, is ten times
the panel estimate $\val{errMesh}$, the larger of the two estimates, fixed
before the comparison.

The largest difference is the reference's ordinate, $\val{cmpRefYDiff}$, then
the empty board's time, $\val{cmpBareTDiff}$, its abscissa,
$\val{cmpBareXDiff}$, and the reference's time, $\val{cmpRefTDiff}$. All are
far below the panel estimate: halving the edge from $\val{panelCoarse}$ moves
our end states by $\val{errMesh}$, while our rectangles of $\val{panelFine}$
and the game's graded triangles land within $\val{cmpRefYDiff}$ of each other.
The estimate is the change from the coarser of our two meshes, so it bounds
the error of the finer one generously. The reference flight is the sensitive
one: it swings round the placed charge, and the field it meets there sets its
direction when it leaves. The step estimate $\val{errStep}$ is below all of
them. The outcomes agree, and the reference's abscissa and the empty board's
ordinate agree exactly, since each is a face the flight meets.

\section{Other solutions}\label{sec:others}

The empty board, the second placement in \code{placements.json}, does not solve
the level. Pushed down by the fixed charge and the glass, the ion leaves
through the \val{BareFace} edge of the world at $(\val{BareX}, \val{BareY})$,
at $t = \val{BareT}$, at $\val{BareAng}$ degrees. Its clearance of the glass
is $\val{BareDmin}$, its energy drifts by $\val{BareWdev}$, and the run takes
$\val{BareSteps}$ steps.

The reference charge without the glass, a run of our own and not a row of the
comparison: the ion ends as \val{HitOutcome} at
$(\val{HitX}, \val{HitY})$, at $t = \val{HitT}$. The glass's extra push sends
the ion down sooner and past the placed charge on the side that turns it up
into the detector; without that push the ion falls onto the charge. A
solution must therefore reckon with the glass, and with the glass's push, not
with a weakened field.

There is no slider, and a panel solve per placement is too slow for a map of
every charge.

\section{What the game computes}\label{sec:engine}

The engine solves the same condition (\physref{2.7}, \emph{Dielectrics}): the
glass's surface cut into flat triangles of constant density, graded towards the
edges, the condition $\sigma = (\lambda/2\pi)\bar E_n$ averaged over each
triangle by a six-point rule of degree four, and each body's neutrality held by
a uniform normal field $\gamma$, a bordered dense system factorized by LU.
Only the upper half is meshed, and the mirror is applied exactly. Preview
flights use panels of half a cell and verification panels of 0.35 of a cell
(\physref{2.7}); its cube polarizability is within about one per cent of the
literature at the verification's resolution.

This document holds the same condition, \eqref{eq:system}, on equal rectangles
with the integral \xeqref{deflection_plates}{eq:rectangle}, averaged at four
Gauss points, with the same neutrality constraint. The two discretizations are
different, so the agreement asked of them is the tolerance $\val{tolFlight}$,
ten times our panel estimate; they agree far more closely
(\S\ref{sec:comparison}). Neither includes the ion's own image (bound below
$\val{RefImage}$ here), and the obstacle is the box plus the contact shell
$\val{contact}$ on both sides.

The flight itself, in the engine, is the adaptive eighth-order Runge--Kutta
method of \physref{5} and \physref{6}, verified as in \physref{7}; once the
field is fixed, a verified end state is good to about $10^{-10}$, and the
difference from ours is the field's, the panels'. Radiation is integrated and
not applied (\xsecref{foundations}{sec:radiation}, \physref{8}).

\section{Further reading}\label{sec:reading}

Dielectrics and their boundary conditions: \Jackson{§4.3}; the image solution
for a point charge and a dielectric half-space, the first step of
\S\ref{sec:slab}: \Jackson{§4.4}. Jackson's §4.3 also states the reduction of a
charge's field by $\varepsilon$ in a medium that fills all space, which is not
this level's case. The expansion of a point charge's potential in Bessel
functions: \Jackson{Ch.~3, cylindrical coordinates}. The motion, the energy, the step and the location of an
event are in the Foundations (\xsecref{foundations}{sec:numerics}); the
rectangle integral and the flight's events are derived in the deflection
plates' document (\xsecref{deflection_plates}{sec:panel},
\xsecref{deflection_plates}{sec:flight}). The engine's dielectrics are
\physref{2.7}.

\levelfileappendix

\end{document}
